\section{An independent proof of source equivalence}
\label{sec:independent-source-bisimulation}

Let $P_0$ denote the complete ordered source list with SHA-256 digest
\begin{quote}\small\ttfamily
74b1e7f73cc3ffe7761e36ab2d57a26e0ba34d3ac624f2472799e6ac41d5a51d,
\end{quote}
and let $P_1$ denote the reparenthesized source list with digest
\begin{quote}\small\ttfamily
7fd152d1f5371d2a561d4ad1870d241683b010ce7b7a27399b95363a4537cfb6.
\end{quote}
Each has $m=6455$ instruction occurrences, listed from left to right in
its full binary decoder tree.  Write $p_i^{(a)}$ for the codeword of
occurrence $i$ in $P_a$, where $0\leq i<m$ and $a\in\{0,1\}$.
The operation and operand at occurrence $i$ are exactly the same in
the two lists.  There are $1931$ decrement occurrences, $2596$
increments, $128$ explicit breaks, $1799$ continues, and one halt.
The maximum codeword lengths are $51$ and $48$, respectively;
$6250$ codewords change.

\subsection{Pointer conventions and a finite semantic quotient}

Pointers are finite binary words, read from their first bit, and
implicitly followed by zeroes.  On decoding a pointer $u$ to a leaf
of depth $d$, the actual RegisterPath semantics consumes
\[
  \operatorname{consume}(u,d)=(u0^d)\mathbin{\upharpoonright}d.
\]
Thus the word on which the instruction operates is precisely its
leaf codeword, even if the incoming pointer is shorter or has an
unused suffix.  For a finite word $p$, let $I(p)$ be its fixed-length
binary increment.  In particular, an all-ones word of length $d$
increments to $0^d$, and $I(\epsilon)=\epsilon$.  Let $S_b(k,p)$
delete the suffix through the $k$th occurrence of bit $b$, counting
from the right; if fewer than $k$ such bits occur, the result is
$\epsilon$.

The successor pointers are $I(p)$ for an increment or positive
decrement, $I(S_0(1,p))$ for a zero decrement,
$I(S_0(k,p))$ for an explicit break of depth $k$, and
$S_1(k,p)$ for a continue of depth $k$.  The halt occurrence has
no successor.  These are the literal definitions in the retained
RegisterPath source, with no replacement of break or continue by
an informal structured-language convention.

For each source, decode every successor pointer with implicit
zero padding and record only the reached instruction index.
This gives a finite control table with $6455$ rows and $8385$
nonhalting successor entries.  The extra $1931$ entries arise by
checking both decrement outcomes.  The one halt row is checked
separately.  Registers do not enter this finite calculation:
the two outcomes of every decrement are retained explicitly.

\paragraph{Independent finite certificate.}
The new checker \texttt{verify\_source\_control.py} compares these
control tables directly.  It does not import the inherited source
checker, replay the rotation certificate, or recurse through a
source tree.  Instead, for the maximum depth $D$ of one source,
it interprets a codeword $p$ as the dyadic integer interval
\[
  J_D(p)=
  [\operatorname{val}(p)2^{D-|p|},
   (\operatorname{val}(p)+1)2^{D-|p|}).
\]
The checker verifies that the ordered intervals partition
$[0,2^D)$ with neither gaps nor overlaps.  This proves both
prefix-freeness and completeness.  A pointer $u$ is decoded by
locating the integer $\operatorname{val}(u)2^{D-|u|}$ in this
partition.  Bit shifts and fixed-width integer addition implement
the pointer operations above.  The checker compares every
operation, operand, guarded successor index, and halt flag,
and checks that both empty pointers decode to index zero.
It pins both input digests and refuses to run with Python
assertions disabled.

The checker also establishes a property needed when connecting
this source argument to exact binary-tape runtime frames.  For
every generated successor $u$, its dyadic integer is exactly the
\emph{left endpoint} of the interval of the decoded leaf $p_i$.
Equivalently,
\[
             u0^\omega=p_i0^\omega.
\]
Consequently replacing the successor by its decoded leaf does
not discard a nonzero unused bit or insert a new one.  This
handles shorter pointers, zero-only unused tails, and all-ones
overflow.  The empty initial pointer likewise agrees with the
first leaf after infinite zero padding.  The check of this
property is separate from equality of instruction indices.

\paragraph{Source equivalence theorem.}
For every register function $R:\mathbb N\to\mathbb N$,
\[
 \operatorname{RegisterHalts}(P_1,(\epsilon,R))
 \quad\longleftrightarrow\quad
 \operatorname{RegisterHalts}(P_0,(\epsilon,R)).
\]
The equivalence requires no zero-register hypothesis and no
bound on any register value.

\paragraph{Proof.}
Project a source configuration to its decoded instruction index
and its entire register function.  Decoding always terminates
at a leaf: both source trees are finite and complete.  In the
pinned Lean realization of $P_0$, all reached DAG children have
smaller indices, its maximum leaf depth is $51$, and the supplied
decoder fuel is $204$.  Thus fuel exhaustion cannot be confused
with the literal halt occurrence.  The source $P_1$ can be
represented by its displayed finite tree with any fuel exceeding
its maximum depth $48$.

The two projections start at the same index zero and the same
register function.  Assume the projections agree at one step.
The occurrence labels agree.  An increment updates the same
register on both sides.  A decrement tests the same register,
therefore selects the same zero/nonzero guard, and applies the
same update.  Breaks and continues do not change registers.
The finite certificate proves that the selected successor
indices agree in every case.  Finally, the unique literal halt
occurrence agrees.  Induction on the number of source steps
therefore gives a bisimulation with exactly the same sequence
of instruction indices, register functions, and halt flags.
This proves the equivalence for finite halting runs, and also
shows that an infinite source execution on either side has an
infinite counterpart.  The registers above index eight are
unchanged throughout; no assumption about them was used.
\hfill$\square$

This proof concerns source instructions.  It does not assert
equal numbers of binary Turing-machine steps: the two decoders
have different codewords and can take different amounts of
time.  Positive whole-tape runtime segment theorems supply the
separate binary connection.

\subsection{The exact inherited consistency calibration}

The calibrated initial register function is
\[
 R_*=(3,0,0,0,3,0,0,0,0),\qquad R_*(i)=0\quad(i\geq9),
\]
and its pointer is empty.  This is the literal value of
\path{LogicalSeed246.initial}, obtained from
\path{L2.enc}\ applied to the record
$(\mathrm{np},\mathrm{tw},\mathrm{ws},\mathrm{ac},
\mathrm{p1},\mathrm{p2},\mathrm{s1},\mathrm{s2},\mathrm{pl})
=(3,0,0,0,3,0,0,0,0)$.
It is not an arbitrary reinterpretation of the binary initializer.
The exact startup/runtime argument must still identify the binary
tape with this logical register function and the empty pointer.

At the pinned ProveIt commit
\texttt{13ec39ed43148b1807f51ed6e6528e1c84659562}, the retained declaration
is \path{CalibratedRotation240.halts_iff_zfc_refutation}.
Write $P_0$ for \path{SourceRotation240.program} and $(\epsilon,R_*)$
for \path{LogicalSeed246.initial}. The declaration states
\[
 \operatorname{RegisterHalts}(P_0,(\epsilon,R_*))
 \quad\longleftrightarrow\quad
 \operatorname{BProv}(\mathrm{ZFCax\_s},[],\bot).
\]
The identity of its program with the complete source $P_0$ is
checked separately by extracting the actual DAG lookup definitions,
including the overrides at rows $5$, $6$, and $169$.  Storage
contains $203$ rows; $202$ internal rows are reachable, and the
unfolding contains exactly the pinned $6455$ leaves.

The declaration composes explicit inherited interfaces.  First,
\path{SourceRotation240.program_halts_iff} identifies the new
source with \path{ProgramRotation244.program} for every register
function.  Then \path{SourceRotation244.program_halts_iff}
and \path{BisimRouting245.program_halts_iff} identify that
source with \path{Prog254.program}, again for every register
function.  At $R_*$,
\path{LogicalSeed246.halts_iff_native_refutation} uses the
existing \path{Bisim254.prog254_halts_iff} and
\path{DirtyProgram254.six_nonformula_path_iff}
interfaces to obtain a refutation of the native first-order
calculus.  The two remaining comparisons are
\begin{align*}
 \operatorname{BProv}(\mathrm{NativeAx},[],\bot)
  &\longleftrightarrow
   \operatorname{BProv}(\mathrm{ZFax\_s},[],\bot),\\
 \operatorname{BProv}(\mathrm{ZFax\_s},[],\bot)
  &\longleftrightarrow
   \operatorname{BProv}(\mathrm{ZFCax\_s},[],\bot),
\end{align*}
The first is provided by \path{WFTheory.native_refutation_iff_zf}.
The second is provided by
\path{ConstructibleChoice.zf_refutation_iff_zfc_refutation}.

The first comparison uses the internally well-founded part of the
CP/GIC quotient to build a model of ZF and uses first-order soundness
and completeness to compare finite refutations.  The second
constructs the constructible submodel of a model of ZF, verifies
Choice there, and again applies soundness and completeness.  In
particular, the actual comparison includes Foundation and Power
on the ZF side and the Choice axiom on the ZFC side.  These
statements are more specific than a bare informal claim that two
theories are equiconsistent.

Finally, $\operatorname{BProv}$ asserts the existence of a finite
list of theory axioms and a finite natural-deduction derivation
from that list, expressed by the inductive relation
$\operatorname{Prov}$.  Encoding finite derivation trees by natural numbers,
and decoding such encodings, identifies this endpoint with
$\exists p\,\operatorname{Proof}_{\mathrm{ZFC}}(p,\bot)$
for the retained first-order proof convention.  Ordinary changes
between standard effective proof encodings are finite syntactic
translations.  Negation gives the corresponding consistency
statement.

\paragraph{Internalization and scope of this audit.}
The new source theorem uses finite arithmetic checks and induction
on source steps, so its ordinary mathematical proof is formalizable
in ZFC.  The inherited model-theoretic comparisons use finite
syntax, first-order models, quotients, and definable submodels;
their ordinary arguments also fit within ZFC.  This observation
does not rely on assuming a general soundness theorem for all of
Lean's type theory inside ZFC.  The present work has inspected
the cited declaration statements and their explicit composition,
and freshly verified the finite source certificate.  It has not
rebuilt or independently re-proved the complete ancestor Lean
library, nor produced a new serialized proof for the entire
binary endpoint. This inherited calibration is the logical
dependency in the final ZFC biconditional. No consistency assumption is needed to prove the
biconditional; actual nonhalting follows under
$\operatorname{Con}(\mathrm{ZFC})$.
